%%%%%%%%%%%%%%%%%%%%%%%%%%%%%%%%%%%%%%%%%%%%%%%%%%%%%%%%%%%%%%%%%%%%%%%%%%%%%%%%
% abstract.tex
%
% Modelo de arquivo para uso com a classe uvvTeX2, para a formatação de
% trabalhos acadêmicos na Universidade Vila Velha (UVV) (https://www.uvv.br).
%
% Para maiores informações, visite:
%    https://github.com/uvv-computacao/uvvtex2
%
% Neste arquivo você deve escrever o abstract de sua monografia, ou seja, deve
% escrever o resumo em INGLÊS. O resumo deve ser objetivo e listar os aspectos
% mais importantes de seu trabalho. Ao final do abstract, as keywords que você
% definiu no arquivo principal da monografia serão inseridas automaticamente.
%%%%%%%%%%%%%%%%%%%%%%%%%%%%%%%%%%%%%%%%%%%%%%%%%%%%%%%%%%%%%%%%%%%%%%%%%%%%%%%%

% Não altere as linhas a seguir:
\setlength{\absparsep}{18pt}
\begin{resumo}[Abstract]
\begin{otherlanguage*}{english}

% Começe a escrever o abstract aqui:
This study presents the analysis and modeling of a pedagogical workflow
automation solution for the Educa EMES system, used by the Escola da
Magistratura do Espírito Santo (EMES). The research starts from the existence
of an integrated platform already in production, while identifying gaps related
to the persistence of manual interventions in stages such as opening and
closing enrollments, participant confirmation, approval, certification, and
report generation. Methodologically, the study is characterized as applied
research with a qualitative approach, based on a case study of a real system
used in an institutional context. The current pedagogical workflow, the
system's technological structure and data model, the main operational
bottlenecks, and the requirements needed to expand automation are analyzed. As
a result, the work proposes a model based on configurable steps, eligibility
analysis, pending issue identification, state traceability, and preservation of
manual intervention when necessary, guiding the development and evaluation
stages of the solution.
 
% Não altere as linhas a seguir:
\textbf{Keywords}: \imprimirkeywords.
\end{otherlanguage*}
\end{resumo}
